% limitations.tex, candid; each item ties to a logged result or a known gap.
\section{Limitations}
\label{sec:limitations}

We state the limitations plainly; several bound the strength of the present claims and motivate the
GPU-scale replication the pipeline is built for.

\begin{itemize}
 \item \textbf{Single, smallest model.} All results are on Pythia-$160$M (CPU). Memorization grows
 log-linearly with model scale~\cite{carlini2023quantifying}, so both the membership signal and the
 extraction outcome are expected to be stronger at $1.4$B--$12$B. The present numbers are
 \emph{preliminary}; we have built every analysis so the larger-model run is a one-line
 configuration change.
 \item \textbf{Chance-level membership separation.} On the confound-clean Pile train-vs-val split,
 membership AUC is at chance ($0.45$--$0.49$) at $160$M, consistent with~\cite{duan2024mia}. The
 divergence result is therefore established in a regime where the membership signal is itself weak;
 whether the calibrated detectors gain \emph{independent} leakage-predictive value once membership
 separation becomes non-trivial at scale is an open question our design is poised to answer.
 \item \textbf{Near-degenerate extraction outcome.} Extractable memorization at $160$M is rare
 ($3/300$ items fully extracted; mean fractional extraction $0.037$), so the correlation analysis
 leans on a small high-extraction tail. We mitigate with rank statistics, bootstrap CIs, and a
 zero-robust Kendall check, but a less zero-inflated outcome at scale would sharpen all estimates.
 \item \textbf{PII not yet demonstrated.} On the Enron-in-Pile subset we observed \emph{zero}
 verbatim PII leakage at $160$M ($8/36$ documents contained PII in the held suffix; none were
 regurgitated). The PII limb of the threat model is thus a designed capability with a null result at
 this scale, not a demonstrated leak; we report it as such and do not claim PII exposure.
 \item \textbf{Benchmark-level test underpowered.} The Oren permutation/exchangeability test is run
 only at sanity scale on $160$M; membership-based, it is underpowered here and is flagged as
 GPU-gated rather than used to draw contamination conclusions.
 \item \textbf{$n$-gram contamination is a lower bound.} Our model-free $n$-gram overlap uses a
 public \emph{sample} of the Pile as the reference index, so measured benchmark$\leftrightarrow$Pile
 overlap underestimates the true overlap against the full corpus.
 \item \textbf{Observational, members-only correlation.} The headline analysis correlates detector
 scores with extraction across known members; it is observational, not interventional. We address
 the most important confound (loss) by pre-registered partial correlation and mediation, and the
 obvious alternatives (frequency, duplication, non-linearity, distribution shift) by explicit
 controls, but residual confounding cannot be excluded.
 \item \textbf{Collinearity of detectors with loss.} The calibrated detectors are deterministic
 transforms of the same per-token log-probabilities as loss and are empirically collinear with it
 (Spearman $0.74$--$0.90$; VIF up to $6.2$ for Min-K\%). Consequently we interpret the negative
 partial/direct terms as possible \emph{suppression artifacts} of near-collinearity and claim only
 the conservative ``no positive residual'' result; we do not assert the detectors inversely predict
 leakage.
 \item \textbf{Construct validity of the leakage proxy.} The outcome (greedy prefix-continuation
 extraction over the held suffix) is itself likelihood-related, so part of the loss$\leftrightarrow$
 extraction association is mechanical/definitional. Our control removes the loss component, but a
 decisive separation would compute prefix-only loss against extraction; we flag this as a known
 construct-validity limitation rather than claiming the two are independent by construction.
 \item \textbf{Selection and aggregation.} Members are drawn from a non-uniform public Pile sample
 (\texttt{pile-10k}), so member-selection bias is possible; and the pooled correlation aggregates
 domains whose effects flip sign (Section~\ref{sec:res-headline}), so the pooled $\rho$ should be
 read as a domain-mixture, not a homogeneous effect.
 \item \textbf{Linearity (now addressed).} An earlier version controlled for loss only linearly; we
 added a cubic-residual and decile-stratified non-linear control, under which no positive
 independent signal revives. We note it here because it was a live threat to the claim until tested.
\end{itemize}
